% !PRINTED 129
\weberpage{129}
\begin{center}
\S.\,32.

\gesperrt{Satz von Frobenius.}
\end{center}

\gesperrt{Frobenius} hat einen Satz aufgestellt, der als ein Gegenstück zum vierten Sylow'schen betrachtet werden kann, der sich so aussprechen lässt\footnote{\gesperrt{Frobenius}, Sitzungsberichte der Berl. Akademie, 4. Mai 1893.}:

V. \gesperrt{Jede Gruppe, deren Grad ein Product von lauter verschiedenen Primzahlen ist, ist metacyklisch.}

Wir setzen also voraus, dass der Grad $n$ einer Gruppe $P$ \gesperrt{durch keine Quadratzahl theilbar} sei. Ist $t$ der grösste Primtheiler von $n$, so giebt es nach dem \gesperrt{Cauchy-Sylow}'schen Satze (\S.\,29) einen Theiler $T$ vom Grade $t$ von $P$, der, weil $t$ eine Primzahl ist, eine cyklische Gruppe ist, also aus den Elementen

\[
T = 1, \; a, \; a^2, \; \ldots \; a^{t-1}
\]

besteht. Wenn sich nun nachweisen liesse, dass unter den gemachten Voraussetzungen $T$ ein Normaltheiler von $P$ ist, so könnte man dieselbe Schlussweise auf die Gruppe $P/T$, deren Grad $n:t$ ist, anwenden, und würde, wenn $t_1$ die zweitgrösste in $n$ aufgehende Primzahl ist, auf einen Normaltheiler $N$ der Gruppe $P/T$ vom Grade $t_1$ schliessen. Daraus würde dann aber folgen (\S.\,6, 2.), dass $P$ einen Normaltheiler $T_1$ vom Grade $t\,t_1$ hat, von dem $T$ selbst wieder Normaltheiler ist. Indem man in gleicher Weise die Gruppe $P/T_1$ betrachtet u. s. f., ergiebt sich eine Compositionsreihe von $P$:

\[
P, \; T_k, \; T_{k-1}, \; \ldots \; T_1, \; T, \; 1,
\]

in der die Indexreihe aus den der Grösse nach aufsteigend geordneten Primfactoren von $n$ besteht, und $P$ ergiebt sich als metacyklisch. Alles kommt also darauf an, nachzuweisen, dass $T$ ein \gesperrt{Normaltheiler} von $P$ ist.

Dieser Satz ist leichter zu beweisen, wenn man ihn als speciellen Fall eines allgemeineren betrachtet, als wenn man ihn direct angreift. Dieser allgemeinere Satz lautet unter den über $P$ und $n$ gemachten Voraussetzungen:

$\alpha$) \gesperrt{Ist $n = \mu\nu$ in zwei Factoren zerlegt und ist jeder Primtheiler von $\nu$ grösser als jeder Prim-}
%CONT>>
